% answers_short.tex —— 2007 年·普通高等学校招生考试·福建卷·理科 · 简短答案速查
% 由 tools/make_short_answers.py 自 answers.tex 的 \gksolhead/\gkanswer 提取，
% 并与 2007/2007_数学真题_福建理_答案要点.md 逐题比对后生成，勿手改（改源 MD 后重跑）。

\begin{gknotice}
\textbf{使用说明}：本节只给最终答案便于快速核对，完整推导见解析版。答案为编写者独立推导，\textbf{非官方标准答案}。
\end{gknotice}

\gkpart{选择题}
\begin{center}\begin{tabular}{ccccccccccccc}
\toprule
题号 & 1 & 2 & 3 & 4 & 5 & 6 & 7 & 8 & 9 & 10 & 11 & 12 \\
\midrule
答案 & D & B & C & B & A & A & C & D & D & B & B & D \\
\bottomrule
\end{tabular}\end{center}

\gkpart{填空题}
\begin{itemize}[leftmargin=2.6em]
  \item[\textbf{13.}] $[-5,7]$。
  \item[\textbf{14.}] $\sqrt2-1$。
  \item[\textbf{15.}] $\dfrac23$。
  \item[\textbf{16.}] 例如：三角形的全等关系、三角形的相似关系、整数的同余关系。
\end{itemize}

\gkpart{解答题}
\begin{itemize}[leftmargin=2.6em]
  \item[\textbf{17.}] （1）$C=135^\circ$；（2）$BC=\sqrt2$。
  \item[\textbf{18.}] （1）证明见解析；（2）$\arccos\dfrac{\sqrt6}4$；（3）$\dfrac{\sqrt2}2$。
  \item[\textbf{19.}] （1）$L=(x-3-a)(12-x)^2$（$9\leqslant x\leqslant11$）；（2）$3\leqslant a\leqslant\dfrac92$ 时 $x=9$、$Q=54-9a$；$\dfrac92\leqslant a\leqslant5$ 时 $x=6+\dfrac{2a}3$、$Q=\dfrac{4(9-a)^3}{27}$。
  \item[\textbf{20.}] （1）$y^2=4x$；（2）$0$。
  \item[\textbf{21.}] （1）$a_n=2n-1+\sqrt2$，$S_n=n(n+\sqrt2)$；（2）证明见解析。
  \item[\textbf{22.}] （1）减区间 $(-\infty,1)$，增区间 $(1,+\infty)$；（2）$0<k<\mathrm e$；（3）证明见解析。
\end{itemize}
